\centerline{\bf \Huge Разнобой по многочленам}

\begin{flushright}
        \scriptsize{Эпизод 7. Созданные человеком}
\end{flushright}

\problem{1}
Дан многочлен $P(x)$. Докажите, что для некоторого многочлена $Q(x) \neq x$ многочлен $P(Q(x))$ делится на $P(x)$.

\problem{2}
Дано вещественное число $a$. Многочлен $P(x)$ с вещественными коэффициентами таков, что $P(x) \equiv P(a-x)$. Докажите, что $P(x) \equiv Q\left(\left(x-\frac{a}{2}\right)^2\right)$ для некоторого многочлена $Q(x)$ с вещественными коэффициентами.

\problem{3}
Пусть $P(x)$ и $Q(x)$ - многочлены с действительными коэффициентами. Известно, что многочлен $P\left(x^3\right)+x Q\left(x^3\right)$ делится на $x^3-1$. Докажите, что $P(x)$ и $Q(x)$ оба делятся на $x-1$.

\problem{4}
Многочлен $P(x) \in \mathbb{Q}[x]$ неприводим над $\mathbb{Q}$. Докажите, что он не может иметь кратный комплексный корень.

\problem{5}
(a) Многочлены $P(x)$ и $Q(x)$ с целыми коэффициентами таковы, что $P(n)$ делится на $Q(n)$ при бесконечном количестве целых $n$. Докажите, что $P(x)$ делится на $Q(x)$ (то есть в частном получается многочлен с рациональными коэффициентами).
(b) Многочлены $P(x)$ и $Q(x)$ с целыми коэффициентами взаимно просты как многочлены над $\mathbb{Q}$. Докажите, что существует некоторая константа $C$ такая, что для всех целых $n$ выполнено $(P(n), Q(n)) \leqslant C$.

\problem{6}
Существуют ли многочлены с вещественными коэффициентами $P(x)$ и $Q(x)$ такие, что

$$
P(x)^2+Q(x)^3=x ?
$$